\chapter{Phát triển hệ thống}

\section{Phần cứng}

Hệ thống được xây dựng dựa trên vi điều khiển ESP32, module đo ECG AD8232,
module thẻ nhớ SD và các ngoại vi hỗ trợ quá trình ghi nhận dữ liệu. Module
AD8232 đảm nhiệm vai trò khuếch đại và tiền xử lý tín hiệu điện tim thu được từ
các điện cực. Ba điện cực chính gồm RA, LA và RL được gắn lên cơ thể người đo;
trong đó RA và LA là hai đầu vào vi sai của tín hiệu ECG, còn RL đóng vai trò
điện cực tham chiếu.

Tín hiệu ECG sau tiền xử lý được lấy ra tại chân \texttt{OUT} của AD8232 và đưa
vào chân ADC của ESP32. Ngoài ra, hai chân \texttt{LO+} và \texttt{LO-} của
AD8232 có thể được dùng để nhận biết tình trạng rời điện cực. Tuy nhiên, trong
thực tế hai chân này có thể xuất hiện xung hoặc trạng thái không ổn định khi
toàn bộ điện cực chưa được gắn, do các đầu vào của mạch bị floating. Vì vậy,
trạng thái điện cực không nên được quyết định bằng một lần đọc GPIO duy nhất mà
cần kết hợp lọc theo cửa sổ thời gian hoặc kiểm tra thêm biên độ tín hiệu ADC.

Chân \texttt{SDN} của AD8232 là chân shutdown. Khi \texttt{SDN} ở mức cao,
module hoạt động bình thường; khi \texttt{SDN} ở mức thấp, module chuyển sang
trạng thái tiết kiệm năng lượng. Trong hệ thống hiện tại, chân này có thể được
kéo lên mức cao hoặc điều khiển bởi một GPIO của ESP32 nếu cần bật/tắt module
bằng phần mềm.

Module thẻ nhớ SD được kết nối với ESP32 thông qua giao tiếp SPI để lưu dữ liệu
ECG dưới dạng file CSV. Ngoài các chân SPI, hệ thống sử dụng thêm chân
\texttt{GPIO35} để nhận tín hiệu DET của module thẻ SD. Khi trạng thái thẻ SD
thay đổi, chân này tạo ngắt để firmware đưa hệ thống về trạng thái khởi tạo lại.
Một LED trạng thái được nối tại \texttt{GPIO21}; LED được bật khi hệ thống đang
ở trạng thái ghi dữ liệu và tắt ở các trạng thái còn lại.

\begin{figure}[H]
    \centering
    \includegraphics[width=0.9\linewidth]{figures/Heart_Rate_Monitor_wiring.png}
    \caption{Sơ đồ đi dây của các module trong hệ thống}
    \label{fig:hardware_wiring}
\end{figure}

\section{Thiết kế hệ thống}

Phần mềm được thiết kế theo hướng tách biệt giữa khối quản lý trạng thái và
khối dịch vụ xử lý ECG. File \texttt{main.c} quản lý state machine chính của hệ
thống, bao gồm các trạng thái \texttt{INIT}, \texttt{IDLE} và \texttt{RECORD}.
Trong khi đó, file \texttt{ecg\_pipeline.c} chứa các hàm dịch vụ mà state
machine gọi khi cần, ví dụ khởi tạo pipeline, bắt đầu ghi dữ liệu, dừng ghi dữ
liệu, reset pipeline, đọc ADC, xử lý ECG và ghi dữ liệu vào thẻ SD.

Cách tổ chức này giúp state machine không bị trộn lẫn với các hàm xử lý thấp
tầng. Các task dịch vụ không tự ý chuyển trạng thái hệ thống mà chỉ phát sinh sự
kiện lỗi về state machine. State machine là nơi duy nhất quyết định hệ thống sẽ
ở trạng thái nào tiếp theo.

\subsection{Thiết kế state machine}

State machine hiện tại gồm ba trạng thái chính:

\begin{itemize}
    \item \textbf{INIT}: trạng thái khởi tạo và phục hồi hệ thống.
    \item \textbf{IDLE}: trạng thái sẵn sàng, chờ người dùng bắt đầu ghi dữ liệu.
    \item \textbf{RECORD}: trạng thái đang lấy mẫu, xử lý và lưu dữ liệu ECG.
\end{itemize}

\begin{figure}[H]
    \centering
    \includegraphics[width=0.95\linewidth]{figures/State_machine.png}
    \caption{Sơ đồ chuyển trạng thái của hệ thống}
    \label{fig:state_machine}
\end{figure}

Trong trạng thái \texttt{INIT}, hệ thống thực hiện các thao tác phục hồi an
toàn. Firmware dừng quá trình ghi dữ liệu nếu đang hoạt động, dừng timer lấy
mẫu, hủy khởi tạo ADC nếu cần, đóng hoặc unmount thẻ SD nếu cần và tắt LED báo
ghi dữ liệu tại \texttt{GPIO21}. Sau đó hệ thống thực hiện kiểm tra sức khỏe
cơ bản. Nếu các điều kiện khởi tạo đạt yêu cầu, hệ thống chuyển sang trạng thái
\texttt{IDLE}; nếu không, hệ thống tiếp tục ở lại \texttt{INIT} và thử kiểm tra
lại sau một khoảng thời gian.

Trong trạng thái \texttt{IDLE}, hệ thống đã sẵn sàng nhưng chưa ghi dữ liệu.
Timer ADC chưa được chạy, ADC chưa lấy mẫu và LED \texttt{GPIO21} ở mức thấp.
Ở trạng thái này, firmware thực hiện health check định kỳ để kiểm tra các điều
kiện cơ bản như heap, stack, trạng thái thẻ SD và các lỗi hệ thống. Khi người
dùng nhấn nút record tại \texttt{GPIO0}, state machine gọi hàm
\texttt{ecg\_pipeline\_enter\_recording()}. Nếu mount thẻ SD và khởi động các
dịch vụ ghi thành công, hệ thống chuyển sang \texttt{RECORD}. Nếu có lỗi, hệ
thống quay về \texttt{INIT}.

Trong trạng thái \texttt{RECORD}, hệ thống mount thẻ SD, khởi tạo ADC, reset
bộ xử lý tín hiệu, bật timer lấy mẫu và bắt đầu ghi dữ liệu. Mỗi khi timer tạo
sự kiện lấy mẫu, task xử lý ECG đọc ADC, xử lý tín hiệu, phát hiện đỉnh R, tính
BPM và gửi mẫu dữ liệu sang task ghi SD. Dữ liệu được ghi xuống thẻ SD theo
dạng CSV. Nếu có BLE subscriber, hệ thống cũng có thể gửi dữ liệu ECG gần thời
gian thực qua BLE. Khi hệ thống ở trạng thái này, LED tại \texttt{GPIO21} được
kéo lên mức cao để báo đang ghi dữ liệu.

\begin{table}[H]
    \centering
    \caption{Các trạng thái chính của hệ thống}
    \label{tab:system_states}
    \begin{tabular}{|p{3cm}|p{10cm}|}
        \hline
        \textbf{Trạng thái} & \textbf{Chức năng} \\
        \hline
        \texttt{INIT} &
        Dừng ghi, dừng timer, hủy khởi tạo ADC nếu cần, reset/close SD nếu cần,
        tắt LED recording và kiểm tra điều kiện phục hồi hệ thống. \\
        \hline
        \texttt{IDLE} &
        Hệ thống sẵn sàng, chưa ghi dữ liệu, LED recording tắt, thực hiện health
        check định kỳ và chờ nút record. \\
        \hline
        \texttt{RECORD} &
        Mount SD, lấy mẫu ADC theo timer, xử lý ECG, ghi CSV, gửi BLE nếu có
        subscriber và bật LED \texttt{GPIO21}. \\
        \hline
    \end{tabular}
\end{table}

\begin{table}[H]
    \centering
    \caption{Sự kiện chuyển trạng thái của hệ thống}
    \label{tab:state_transitions}
    \begin{tabular}{|p{4cm}|p{3cm}|p{3cm}|p{4cm}|}
        \hline
        \textbf{Sự kiện} & \textbf{Trạng thái nguồn} & \textbf{Trạng thái đích} & \textbf{Ghi chú} \\
        \hline
        Khởi tạo OK / system ready & \texttt{INIT} & \texttt{IDLE} &
        Health check đạt yêu cầu. \\
        \hline
        Init fail & \texttt{INIT} & \texttt{INIT} &
        Tiếp tục ở lại INIT và thử lại sau. \\
        \hline
        Nhấn nút record + mount SD OK & \texttt{IDLE} & \texttt{RECORD} &
        Nút \texttt{GPIO0} active-low, có debounce phần mềm. \\
        \hline
        Nhấn nút stop & \texttt{RECORD} & \texttt{IDLE} &
        Dừng ghi, dừng timer và tắt LED recording. \\
        \hline
        GPIO35 SD DET interrupt & Mọi trạng thái & \texttt{INIT} &
        Xảy ra khi thẻ SD bị rút/cắm hoặc tín hiệu DET thay đổi. \\
        \hline
        Lỗi mount/ghi SD & \texttt{RECORD} & \texttt{INIT} &
        Pipeline phát sinh sự kiện lỗi SD về state machine. \\
        \hline
        Lỗi ADC / buffer / task & \texttt{RECORD} hoặc \texttt{IDLE} & \texttt{INIT} &
        Lỗi phần mềm được raise thành event để state machine xử lý. \\
        \hline
    \end{tabular}
\end{table}

\subsection{Thiết kế ngắt phần cứng}

Hệ thống sử dụng ngắt phần cứng cho chân DET của thẻ SD. Chân này được kết nối
với \texttt{GPIO35} của ESP32 và được cấu hình ngắt cả hai cạnh
(\textit{any edge}). Khi người dùng rút hoặc cắm thẻ SD, tín hiệu DET thay đổi,
ISR của \texttt{GPIO35} được gọi.

Trong ISR, firmware không thực hiện các thao tác nặng như mount/unmount SD,
dừng ADC hoặc ghi log dài. ISR chỉ đặt cờ sự kiện và đánh thức task quản lý
trạng thái bằng cơ chế notify của FreeRTOS. Sau đó, state task chạy ở ngữ cảnh
task bình thường sẽ xử lý sự kiện này bằng cách chuyển hệ thống về trạng thái
\texttt{INIT}. Thiết kế này giúp ISR ngắn, an toàn và tránh thực hiện các hàm
không phù hợp trong ngữ cảnh ngắt.

Luồng xử lý ngắt phần cứng có thể mô tả như sau:

\begin{verbatim}
GPIO35 SD DET edge
-> ISR sets sd_detect_event_pending = true
-> ISR notifies state task
-> state task handles event
-> enter INIT
-> stop recording services
-> close/unmount SD if needed
-> LED GPIO21 LOW
\end{verbatim}

\subsection{Thiết kế ngắt phần mềm và sự kiện lỗi}

Bên cạnh ngắt phần cứng, hệ thống sử dụng cơ chế sự kiện phần mềm để báo lỗi từ
pipeline về state machine. Các task xử lý ECG và ghi SD không tự đổi trạng thái
hệ thống. Khi xảy ra lỗi, chúng gọi hàm raise event trong pipeline. Callback đã
đăng ký trong \texttt{main.c} sẽ nhận event, đặt cờ pending và notify state
task. State task sau đó quyết định chuyển trạng thái.

Các loại sự kiện lỗi chính gồm:

\begin{itemize}
    \item \texttt{ECG\_PIPELINE\_EVENT\_SD\_FAULT}: lỗi mount SD, ghi header,
    ghi segment hoặc ghi dữ liệu.
    \item \texttt{ECG\_PIPELINE\_EVENT\_ADC\_FAULT}: lỗi đọc ADC.
    \item \texttt{ECG\_PIPELINE\_EVENT\_BUFFER\_FAULT}: ring buffer đầy và rơi
    mẫu quá nhiều.
    \item \texttt{ECG\_PIPELINE\_EVENT\_TASK\_FAULT}: lỗi tạo task hệ thống.
    \item \texttt{ECG\_PIPELINE\_EVENT\_BLE\_FAULT}: lỗi khởi tạo BLE. Lỗi này
    hiện chỉ được log, không bắt buộc kéo hệ thống về \texttt{INIT} vì SD
    pipeline vẫn có thể hoạt động.
\end{itemize}

Luồng xử lý sự kiện phần mềm:

\begin{verbatim}
pipeline/service fault
-> ecg_pipeline_raise_event(...)
-> callback in main.c
-> set pipeline_event_pending bit
-> notify state task
-> state task checks event type
-> severe fault: enter INIT
\end{verbatim}

Cách thiết kế này giúp state machine giữ vai trò trung tâm. Các module dịch vụ
chỉ báo lỗi, còn quyết định phục hồi hoặc chuyển trạng thái được thực hiện ở
một nơi duy nhất.

\subsection{Thiết kế timer lấy mẫu}

Việc lấy mẫu tín hiệu ECG được điều khiển bằng GPTimer của ESP32. Timer được
cấu hình để tạo sự kiện định kỳ tương ứng với tần số lấy mẫu mục tiêu khoảng
360 Hz. Khi timer đến chu kỳ lấy mẫu, callback timer không đọc ADC trực tiếp mà
chỉ phát semaphore cho task xử lý ECG.

Task xử lý ECG chờ semaphore từ timer. Khi nhận được semaphore, task kiểm tra
hệ thống có đang ở trạng thái ghi dữ liệu hay không. Nếu \texttt{record\_active}
đang bật, task đọc ADC, gọi hàm xử lý tín hiệu ECG và gửi mẫu đã xử lý sang
ring buffer để task ghi SD lưu xuống thẻ nhớ.

Luồng timer lấy mẫu:

\begin{verbatim}
GPTimer alarm at 360 Hz
-> timer callback gives adc_timer_semaphore
-> ECG processing task wakes up
-> read ADC
-> process ECG sample
-> send sample to SD ring buffer
\end{verbatim}

Thiết kế này tách biệt phần ngắt timer và phần xử lý nặng. Timer callback chỉ
làm nhiệm vụ đánh thức task, còn việc đọc ADC, lọc tín hiệu và ghi dữ liệu được
thực hiện trong task bình thường.

\section{Xử lý tín hiệu số}

Tín hiệu ECG sau khi được module AD8232 khuếch đại sẽ được ESP32 lấy mẫu thông
qua ADC. Dữ liệu ADC thô được xử lý qua nhiều bước để giảm nhiễu, loại bỏ trôi
nền và phát hiện đỉnh R. Luồng xử lý hiện tại gồm dịch mức ADC, ước lượng
baseline, loại bỏ baseline, lọc số và phát hiện R-peak.

Đầu tiên, giá trị ADC thô được dịch quanh mức giữa ADC để tạo tín hiệu
\texttt{raw\_centered}. Sau đó, firmware sử dụng bộ lọc trung bình trượt để
ước lượng baseline. Tín hiệu sau khi trừ baseline được gọi là
\texttt{corrected}. Tín hiệu này tiếp tục đi qua các bộ lọc high-pass,
stop-band/notch, low-pass và Kalman để tạo ra tín hiệu \texttt{filtered}.

Luồng xử lý tín hiệu hiện tại:

\begin{verbatim}
raw_adc
-> raw_centered = raw_adc - 2048
-> baseline = moving_average(raw_centered)
-> corrected = raw_centered - baseline
-> high-pass filter
-> stop-band / notch filter
-> low-pass filter
-> Kalman filter
-> filtered
-> R-peak detection
-> RR interval / BPM
\end{verbatim}

Các dữ liệu chính được lưu xuống thẻ SD dưới dạng CSV. Header CSV hiện tại gồm:

\begin{verbatim}
timestamp_ms,raw_adc,baseline,corrected,filtered,is_peak,bpm
\end{verbatim}

Trong đó, \texttt{timestamp\_ms} là thời điểm lấy mẫu, \texttt{raw\_adc} là
giá trị ADC thô, \texttt{baseline} là giá trị nền ước lượng,
\texttt{corrected} là tín hiệu sau khi loại bỏ baseline, \texttt{filtered} là
tín hiệu sau các bộ lọc, \texttt{is\_peak} đánh dấu mẫu được phát hiện là đỉnh
R và \texttt{bpm} là nhịp tim ước lượng.

Trong các giai đoạn tiếp theo, hệ thống có thể bổ sung thêm các cột như
\texttt{raw\_centered}, \texttt{qrs\_energy}, \texttt{qrs\_threshold} và
\texttt{quality\_bad} để hỗ trợ debug thuật toán phát hiện R-peak và nhận biết
motion artifact hoặc rời điện cực.

\subsection{Datapath xử lý dữ liệu ECG}

Datapath của hệ thống mô tả đường đi của dữ liệu từ cảm biến ECG đến các khối
xử lý, lưu trữ và truyền thông. Trong thiết kế hiện tại, dữ liệu bắt đầu từ tín
hiệu analog tại chân \texttt{OUT} của module AD8232, sau đó được ESP32 lấy mẫu
bằng ADC. Mỗi mẫu ADC được đưa vào khối xử lý tín hiệu số để loại bỏ trôi nền,
lọc nhiễu và phát hiện R-peak. Kết quả cuối cùng được lưu vào thẻ SD dưới dạng
CSV và có thể gửi qua BLE nếu có thiết bị nhận dữ liệu.

\begin{figure}[H]
    \centering
    \includegraphics[width=0.95\linewidth]{figures/ECG_datapath.png}
    \caption{Datapath xử lý dữ liệu ECG trong hệ thống}
    \label{fig:ecg_datapath}
\end{figure}

Nếu chưa có hình datapath riêng, datapath có thể mô tả bằng lưu đồ text như sau:

\begin{verbatim}
AD8232 electrodes
-> AD8232 analog front-end
-> OUT analog ECG
-> ESP32 ADC
-> raw_adc
-> baseline removal
-> digital filters
-> filtered ECG
-> R-peak detector
-> RR interval / BPM
-> SD CSV storage
-> BLE notification
\end{verbatim}

Trong datapath này, khối phát hiện R-peak được đặt sau các bộ lọc số. Lý do là
tín hiệu ECG thô từ ADC còn chứa trôi nền, nhiễu cơ, nhiễu tiếp xúc điện cực và
nhiễu cao tần. Nếu phát hiện R-peak trực tiếp trên tín hiệu thô, thuật toán dễ
nhận nhầm motion artifact hoặc spike nhiễu thành đỉnh R. Vì vậy, tín hiệu được
lọc và hiệu chỉnh trước, sau đó mới đưa vào thuật toán phát hiện QRS/R-peak.

\subsection{Lưu đồ thuật toán xử lý ECG}

Thuật toán xử lý ECG trong firmware được thiết kế theo hướng xử lý từng mẫu
theo thời gian thực. Timer tạo nhịp lấy mẫu định kỳ, task xử lý ECG đọc ADC và
đưa mẫu vào chuỗi xử lý số. Sau khi tính được tín hiệu \texttt{filtered}, hệ
thống gọi bộ phát hiện R-peak để xác định nhịp tim.

\begin{figure}[H]
    \centering
    \includegraphics[width=0.95\linewidth]{figures/ECG_algorithm_flow.png}
    \caption{Lưu đồ thuật toán xử lý tín hiệu ECG}
    \label{fig:ecg_algorithm_flow}
\end{figure}

Lưu đồ thuật toán có thể mô tả như sau:

\begin{verbatim}
Start RECORD
-> Start ADC timer
-> Wait timer semaphore
-> Read ADC sample
-> Convert raw_adc to raw_centered
-> Estimate baseline by moving average
-> corrected = raw_centered - baseline
-> Apply HP / notch / LP / Kalman filters
-> Run R-peak detector
-> If R-peak:
       update RR interval
       update BPM
       update AFib/RR irregularity features
-> Push sample to SD ring buffer
-> SD task writes CSV batch
-> Repeat until STOP or fault event
\end{verbatim}

Trong trường hợp phát hiện lỗi như lỗi ADC, lỗi ghi SD hoặc ring buffer đầy quá
nhiều lần, pipeline sẽ phát sinh sự kiện lỗi phần mềm. State machine nhận sự
kiện này và chuyển hệ thống về trạng thái \texttt{INIT} để phục hồi.

\subsection{Vị trí của thuật toán Pan-Tompkins trong pipeline}

Thuật toán Pan-Tompkins hoặc các biến thể như Pan-Tompkins++ không nằm ở tầng
đọc ADC, ghi SD hay quản lý trạng thái. Thuật toán này thuộc khối xử lý tín hiệu
số, cụ thể là khối phát hiện QRS/R-peak. Vì vậy, trong kiến trúc firmware, nó
nên được đặt bên trong module xử lý tín hiệu, ví dụ trong file
\texttt{ecg\_signal.c} hoặc tách riêng thành module mới như
\texttt{ecg\_rpeak\_detector.c}.

Vị trí hợp lý của Pan-Tompkins trong datapath:

\begin{verbatim}
raw_adc
-> baseline removal
-> digital filters
-> filtered ECG
-> Pan-Tompkins / Pan-Tompkins light / Two-MA detector
-> R-peak index
-> RR interval
-> BPM
\end{verbatim}

Không nên đặt Pan-Tompkins trong \texttt{main.c}, vì \texttt{main.c} chỉ nên
quản lý state machine và sự kiện hệ thống. Cũng không nên đặt thuật toán trực
tiếp trong driver ADC hoặc driver timer, vì các driver này chỉ chịu trách nhiệm
lấy mẫu và phát nhịp thời gian.

Trong bản thử nghiệm offline, nhóm đã kiểm tra các hướng phát hiện R-peak gồm
Pan-Tompkins++, Pan-Tompkins light viết bằng C và Two Moving Average QRS
detector. Kết quả cho thấy các thuật toán này có thể phát hiện R-peak tốt ở các
đoạn tín hiệu sạch, nhưng khi xuất hiện motion artifact hoặc nhiễu tiếp xúc điện
cực, cần kết hợp thêm các luật kiểm tra chất lượng tín hiệu như refractory time,
kiểm tra biên độ, kiểm tra độ rộng QRS và cờ \texttt{quality\_bad}.

\subsection{Tổng quan thuật toán Pan-Tompkins}

Pan-Tompkins là thuật toán phát hiện QRS kinh điển, phù hợp với hệ thống nhúng
vì có thể triển khai theo từng mẫu và không yêu cầu mô hình học máy. Ý tưởng
chính của thuật toán là làm nổi bật vùng QRS bằng các bước lọc, đạo hàm, bình
phương, tích phân cửa sổ và ngưỡng thích nghi.

Các bước chính của Pan-Tompkins:

\begin{enumerate}
    \item \textbf{Band-pass filtering}: giảm nhiễu trôi nền và nhiễu cao tần,
    giữ lại dải tần chứa năng lượng QRS.
    \item \textbf{Derivative}: tính độ dốc của tín hiệu để làm nổi bật cạnh lên
    và cạnh xuống của QRS.
    \item \textbf{Squaring}: bình phương tín hiệu đạo hàm để biến mọi biên độ
    thành dương và nhấn mạnh các thay đổi nhanh.
    \item \textbf{Moving window integration}: tích phân/trung bình trượt trên
    cửa sổ ngắn để ước lượng năng lượng QRS theo thời gian.
    \item \textbf{Adaptive threshold}: dùng ngưỡng thích nghi để phân biệt vùng
    QRS và nhiễu.
    \item \textbf{Search-back và refractory}: tìm lại nhịp bị bỏ sót và ngăn
    không cho phát hiện hai R-peak quá gần nhau.
\end{enumerate}

Trong hệ thống hiện tại, do tín hiệu đã đi qua các bộ lọc high-pass,
stop-band/notch, low-pass và Kalman, phiên bản triển khai trên MCU có thể dùng
dạng Pan-Tompkins light:

\begin{verbatim}
filtered ECG
-> derivative
-> square
-> moving window integration
-> adaptive threshold
-> remap candidate to max(abs(filtered))
-> refractory check
-> R-peak
\end{verbatim}

Việc remap candidate về \texttt{max(abs(filtered))} trong một cửa sổ ngắn giúp
đưa vị trí phát hiện từ vùng năng lượng QRS về đúng vị trí đỉnh R trên tín hiệu
ECG. Refractory time được dùng để tránh trường hợp thuật toán phát hiện nhiều
đỉnh trong cùng một phức bộ QRS. Với tần số lấy mẫu khoảng 360 Hz, khoảng cách
tối thiểu giữa hai R-peak nên đặt khoảng 200--250 ms thay vì chỉ vài mili-giây.

\subsection{Đề xuất module hóa bộ phát hiện R-peak}

Để firmware dễ mở rộng, bộ phát hiện R-peak nên được tách thành module riêng.
Cấu trúc đề xuất:

\begin{verbatim}
main/
  ecg_signal.c
  ecg_signal.h
  components/
    rpeak/
      ecg_rpeak_detector.c
      ecg_rpeak_detector.h
\end{verbatim}

Module \texttt{ecg\_rpeak\_detector} sẽ cung cấp các hàm:

\begin{verbatim}
void ecg_rpeak_init(uint32_t fs_hz);
void ecg_rpeak_reset(uint32_t now_ms);
ecg_rpeak_result_t ecg_rpeak_process(int16_t filtered, uint32_t now_ms);
\end{verbatim}

Khi đó, \texttt{ecg\_signal.c} chỉ cần gọi detector sau khi đã tạo ra tín hiệu
\texttt{filtered}. Cách này giúp có thể thay đổi thuật toán từ detector hiện
tại sang Pan-Tompkins light, Two Moving Average hoặc fusion detector mà không
ảnh hưởng đến ADC, timer, SD, BLE và state machine.
